\section*{第 \ref{ch:b3:photons-phonons} 章——光子与声子}

\begin{solution}{exo:b3:photons-phonons:1}
(a) 令 $\dd u/\dd\lambda = 0$，得 $x\eu^x/(\eu^x - 1) = 5$，
$x = hc/\lambda k_{\text{B}}T = 4.965$，即 $\lambda_{\max}T =
hc/4.965k_{\text{B}} = \qty{2898}{\micro m.K}$。(b) $0.50$，
$1.07$，$9.35\,\unit{\micro m}$ 与 \qty{1.06}{mm}。(c) 仅太阳的
峰值接近可见光；灯丝峰值位于近红外，仅
以光的形式输出功率的百分之几——这种低效率
使其退役。(d) 人与建筑在十微米处“发光”：
热像仪即为探测第二个峰的相机。
\end{solution}

\begin{solution}{exo:b3:photons-phonons:2}
(a) $\sigma A(T_1^4 - T_2^4) \approx \qty{135}{W}$。(b) 房间的
回流已包含在上述差值之中；衣物引入
温暖的中间表面——辐射是双向的收支账。(c)
$T_{\text{有效}} = [(0.7 \times 1361)/(4\sigma)]^{1/4} =
\qty{255}{K}$。(d) \qty{33}{K} 的额外增温即为温室效应：
红外“保温毯”，
即\cref{pb:b3:harmonic-oscillator:1}所述，并在
\cref{exo:b3:photons-phonons:9}中定量。
\end{solution}

\begin{solution}{exo:b3:photons-phonons:3}
(a) $\num{2.03e7} \times 2.725^3 \approx \qty{4.1e8}{m^{-3}} =
411$ 每 cm$^3$。(b) $1361/\qty{2.2e-19}{J} \approx
\num{6e21}$ 光子每平方米每秒。(c) $3\,\unit{W}/
\qty{3.6e-19}{J} \approx \num{8e18}$ 每秒。(d) 天文
信号每次只到达一个光子——其速率可被计数器分辨；
而加热器每秒 $10^{22}$ 个光子，对任何探测器而言都近似连续。
\end{solution}

\begin{solution}{exo:b3:photons-phonons:4}
(a) 每个模式的 $k_{\text{B}}T$ 乘以模式数 $\nu^2$：
积分后为无穷大——即紫外灾难。(b) 对于 $h\nu \ll
k_{\text{B}}T$。(c) $h\nu/k_{\text{B}}T = 0.025$：处于
经典尾部，此时强度 $\propto T$——因此
以开尔文计的“亮度温度”。(d) $2.8$：峰值
位于普朗克量子折扣全面生效之处。
\end{solution}

\begin{solution}{exo:b3:photons-phonons:5}
(a) 对定向太阳光束，$P \approx \Phi/c =
\qty{4.5}{\micro Pa}$：比大气压强低十一个数量级——
帆之所以可行，只因太空几乎排除了其他一切作用。(b)
$aT^4/3 \approx \qty{1.3e13}{Pa}$：不到气体压强的
千分之一——太阳靠气体压强支撑。(c)
$P_{\text{辐射}}/P_{\text{气}} \propto T^3/n$：质量更大的恒星
燃烧得更烈，$T^4$ 的辐射洪流最终
足以与引力对气体的束缚相抗衡——辐射不稳定性使恒星
质量上限接近 $\sim100\,M_\odot$。(d) $F = P/c = \qty{3.3}{\micro N}$：
牵引光束仍属虚构；光镊利用的是场的\emph{梯度}
作用于微观物体，此时皮牛量级的力占主导
（见\cref{pb:b3:microcanonical-ensemble:1}）。
\end{solution}

\begin{solution}{exo:b3:photons-phonons:6}
(a) $1 + z = 3000/2.725 \approx 1100$。(b) 每个模式的 $\lambda$
都以相同因子伸展，$\nu^3$ 形状映射到自身，
其中 $T \to T/(1+z)$——普朗克曲线不会因膨胀而失去
普朗克形态。(c) $\approx\qty{0.97}{\micro m}$ 的深红峰值
到达时为 \qty{1.06}{mm}。(d) 同一片天空在
\qty{0.5}{\micro m} 处为空，而在毫米处饱和：黑暗
是一种关于波长的表述。
\end{solution}

\begin{solution}{exo:b3:photons-phonons:7}
(a) $\theta_{\text{D}} \approx \qty{500}{K}$，对比拟合值
\qty{343}{K}：量级正确（对纵波与横波速度做恰当平均
后即可弥合差距）。(b) $\approx\qty{2200}{K}$——
与之相符：金刚石是极端情形。(c) 声速低（即软）
和原子质量大两者都使 $\theta_{\text{D}}$ 降低：铅的晶格
几乎立即呈经典行为。(d) $\theta_{\text{D}} \propto
c_{\text{s}}n^{1/3} \propto \sqrt{k/m}$ 的形式，刚而轻
者凝固得晚，软而重者凝固得早。
\end{solution}

\begin{solution}{exo:b3:photons-phonons:8}
(a) 电子贡献：$\gamma \approx \qty{5e-4}{J/(mol.K^2)}$ 给出
\qty{0.5}{mJ/(mol.K)}；晶格贡献：$1944R\dots$——德拜 $T^3$
项给出 $\approx\qty{0.05}{mJ/(mol.K)}$。(b) 电子的贡献
大一个数量级。(c) 在 \qty{3}{K} 附近
（见\cref{exo:b3:quantum-statistics:5}）。(d) 低于交叉温度时
所测 $C/T$ 截距为 $\gamma \propto g(E_{\text{F}})$：
温度计与加热器的实验，读出的是费米面
处的\omterm{prop:b3:schrodinger-three-dimensions:counting}{态密度}。
\end{solution}

\begin{solution}{exo:b3:photons-phonons:9}
(a) 薄层：吸收 $\sigma T_s^4$，发射 $2\sigma T_a^4$（两
面）：$T_s^4 = 2T_a^4$。表面：$S + \sigma T_a^4 = \sigma
T_s^4$。(b) 消去 $T_a$：$S = \tfrac12\sigma T_s^4$，即
$T_s = 2^{1/4}T_{\text{有效}}$。(c) $1.19 \times 255 =
\qty{303}{K}$——高于实际的 \qty{288}{K}，因为
实际保温层是局部的、分层的。(d) 三原子痕量气体的
振动量子：在 $0.5\,\unit{\micro m}$ 处透明，
在 $15\,\unit{\micro m}$ 处不透明——
\cref{pb:b3:harmonic-oscillator:1}中的分子。
\end{solution}

\begin{solution}{exo:b3:photons-phonons:10}
(a) $u = (8\pi k_{\text{B}}^4T^4/h^3c^3)\int x^3/(\eu^x - 1)\dd x
= 8\pi^5k_{\text{B}}^4T^4/15h^3c^3$。(b--c) $\sigma = ac/4 =
2\pi^5k_{\text{B}}^4/15h^3c^2 = \qty{5.67e-8}{W.m^{-2}.K^{-4}}$
——由三个常数导出，恰好吻合。(d) 曲线的\emph{形状}
决定 $h/k_{\text{B}}$（经由 $h\nu/k_{\text{B}}T$），而其
绝对\emph{标度}决定 $k_{\text{B}}^4/h^3$：两个方程，
两个常数——普朗克由同一次拟合读出 $h$ \emph{与} $N_{\text{A}} =
R/k_{\text{B}}$，比佩兰的颗粒实验早数年。
\end{solution}

\begin{solution}{exo:b3:photons-phonons:11}
(a) $B_{12}uN_1 = B_{21}uN_2 + AN_2$。(b) 解出 $u$：
$u = (A/B_{21})/[(B_{12}/B_{21})\eu^{h\nu_0/k_{\text{B}}T} - 1]$;
对所有 $T$ 都与普朗克公式相符，要求 $B_{12} = B_{21}$ 且
$A/B = 8\pi h\nu_0^3/c^3$。(c) 若没有 $A$，则无法
与普朗克定律取得平衡：平衡本身要求受激原子
自发衰变，其速率与受激辐射的速率相关——
爱因斯坦由热力学预言了自发辐射，
比量子力学推导它早了十年。(d) $A \propto
\nu^3$：紫外态在纳秒量级内衰变，
（见\cref{exo:b3:perturbation-theory:8}），因此在短波长下
粒子数反转的代价高得惊人——这正是 X 射线激光的
难题。
\end{solution}

\begin{solution}{exo:b3:photons-phonons:12}
(a) $\kappa = \tfrac13 \times \num{1.8e6} \times \num{1.3e4}
\times \num{3e-7} \approx \qty{2300}{W/(m.K)}$：约为
铜的五倍。(b) 刚而轻意味着声子速度快；纯净、完美、
轻的晶格几乎无物可散射它们——$\ell$ 达到
数千个晶格常数。(c) 即便是
$^{13}$C 原子的质量差异也能可测地散射声子：同位素无序是
真实的热阻，经提纯即可消除。(d) 散热片
位于大功率芯片下方；且嘴唇的热量散失如此之快，
以至于真钻摸上去冷得出奇——珠宝商的
触摸测试，本质上就是一次声子测量。
\end{solution}

\begin{solution}{pb:b3:photons-phonons:1}
\textbf{1.} $\lambda_{\max} = 2898/2.725 \approx \qty{1.06}{mm}$；
$411$ 个光子，每 cm$^3$。
\textbf{2.} $u = aT^4 = \qty{4.2e-14}{J/m^3} \approx
\qty{0.26}{eV/cm^3}$；质量当量 $u/c^2 \approx
\qty{4.6e-31}{kg/m^3}$：约为临界值的 $5\times10^{-5}$。
\textbf{3.} 热谱要求发射体与辐射达到热平衡——
这对任何恒星集合都不可能，
尘埃或等离子体沿不同视线加入亦然；只有一个
\emph{自身曾是}炽热不透明空腔的宇宙才与之相符。
\textbf{4.} $1 + z = 3000/2.725 \approx 1100$。
\textbf{5.} 萨哈平衡转向中性原子：自由电子
消失，汤姆孙散射停止，光子
此后便未受扰动地飞行。
\textbf{6.} 在毫米波段明亮，于
\qty{2.725}{K}时：奥伯斯的天空\emph{确实}炽亮——红移把
这片炽光移出了可见光范围。
\textbf{7.} 每个重子对应 $\num{4.1e8}/0.25 \approx \num{1.6e9}$
个光子。
\textbf{8.} 光子是物质–反物质
湮灭后留下的余烬；重子则是幸存下来的十亿分之一的过剩：
$\eta$ 度量的是原初不对称性——至今
仍未得到解释。
\textbf{9.} $\qty{4.2e-14}{}$ 对比 $\qty{2e-15}{J/m^3}$：
就能量密度而言，这些遗留的光超过了恒星
曾经发出的全部能量。
\textbf{10.} 辐射按 $a^{-4}$ 稀释（数密度 $a^{-3}$，
每个光子的波长拉伸 $a^{-1}$），物质则按 $a^{-3}$ 稀释：早期
辐射占主导；在等效时刻物质超越辐射，使结构
得以增长。
\textbf{11.} $1 + z_{\text{等效}} \approx 3400$。
\textbf{12.} $T_\nu = (4/11)^{1/3} \times 2.725 =
\qty{1.95}{K}$：预言中的宇宙中微子背景，是热物理学中
最大胆、尚未兑现的一张支票。
\textbf{13.} $v = c\,\Delta T/T = \num{3e8} \times \num{1.25e-3}
\approx \qty{370}{km/s}$：太阳系
相对于遗留辐射参考系的运动。
\textbf{14.} 温度（进而是密度与速度）的
涨落，即宇宙在 380\,000 岁时，
是种子的一张照片。
\textbf{15.} 过密区域引力更强，因而增长（
引力不稳定性）；在等效时刻之前，辐射压
将其弹回——增长须等到物质占据主导。
\textbf{16.} 声波（光子–重子流体中的 $c/\sqrt3$）
在复合时刻前所走过的距离——声学视界，
一把埋在天空中的声学标尺。
\textbf{17.} 该角度与平坦宇宙的几何吻合至
约百分之一：宇宙尺度上，平行线依然保持平行。
\textbf{18.} 单极项：炽热的开端。偶极项：我们的
\qty{370}{km/s}。涨落：结构的起源与空间
几何——每一位小数都是一项发现。
\textbf{19.} 它在各个方向上都相同，无论昼夜，
无论夏季冬季，也不受鸽子干扰：没有局域的、银河的或仪器的
效应能如此各向同性——只可能来自天空本身。
\textbf{20.} 每台未调谐的电视机上都有嘶嘶声：大爆炸
一直都在电视上。
\textbf{21.} 热星系团中的电子以康普顿方式使 CMB 光子
频率升高：在交叉频率以下，星系团移除光子
——形成一个深度与距离无关的剪影，
使巡天能在任意红移处找到星系团。
\textbf{22.} 原初引力波会把 CMB 的偏振扭成
“B 模”：相对论的波在
汤姆孙散射的光中被读出——本书的两条线索在此交织。
\textbf{23.} $n_\gamma c/4 \approx \qty{3e16}{m^{-2}s^{-1}}$：
每秒约有 $10^{16}$ 个大爆炸光子穿过你的身体，
而你毫无察觉。
\textbf{24.} 这是一个由\emph{内容}选定的参考系
（辐射各向同性之处），而非由定律选定：相对论
只禁止定律上有特殊地位的参考系——穿过 CMB 运动
是可以探测的，穿过\omterm{thm:b3:relativistic-kinematics:interval}{时空}运动则不可以。
\textbf{25.} 每条视线上都是一条 \qty{2.725}{K} 的普朗克曲线
——每立方厘米 $411$ 个光子，每个重子 $10^9$ 个，
涨落为 $10^{-5}$：统计物理的黑体，被提升
为宇宙学的基础文献。
\end{solution}
